\section{第\ref{sec:acet}章　\nt{ACET}を加える}

アセチルコリンの三つの作用はスライスと生体脳で測定されている。書き込みと読み出しの衝突は本書のモデルで示した。\nt{ACET}が書き込み許可線として働き、予期された不確実性を伝えることは、部分的に実験の裏づけがある仮説である。

{\footnotesize
\setlength{\tabcolsep}{3pt}\renewcommand{\arraystretch}{1.15}
\begin{longtable}{>{\raggedright}p{0.5cm}>{\raggedright}p{4.0cm}>{\raggedright}p{5.6cm}>{\raggedright}p{2.3cm}>{\raggedright\arraybackslash}p{1.7cm}}
\toprule
No. & 主張 & 実験と測定内容 & 出典 & 状態 \\
\midrule\endfirsthead
\toprule
No. & 主張 & 実験と測定内容 & 出典 & 状態 \\
\midrule\endhead
1 & 脳の\nt{ACET}ニューロンは少なく、いくつかの群にまとまり、出力は皮質全体に広がる：ブロードキャストチャネル & ラットでのアセチルコリン合成酵素に対する抗体染色と逆行性標識。皮質、海馬、扁桃体、嗅球、視床は、主として前脳基底部と脳幹にある六つのコンパクトな細胞群からアセチルコリンを受け取る & \cite{Mesulam1983} & 測定済み \\
2 & 皮質のアセチルコリンレベルは覚醒時に高く、徐波睡眠で低い & 自由行動下のネコの皮質と海馬でのマイクロダイアリシスとEEG記録。アセチルコリン放出は徐波睡眠に比べ、静かな覚醒で$\sim100\%$、活動的な覚醒で$\sim175\%$高い。レム睡眠の皮質では覚醒時と同程度 & \cite{Marrosu1995,Hasselmo1999} & 測定済み \\
3 & 信号の意味は受容体が決める：アセチルコリンには速いチャネル型受容体と、細胞内の反応を起動する遅い受容体がある（命題~\ref{prop:receptor}） & 測定のレビュー。興奮性、伝達、可塑性に対するアセチルコリンの作用は、放出部位、受容体サブタイプ（ニコチン性はイオンチャネル、ムスカリン性はGタンパク質経由）、受け手の細胞に依存する & \cite{Picciotto2012} & 測定済み \\
4 & 第一の作用：外部入力にほとんど触れずに内部結合を弱める（\eqref{eq:acetnet}の係数$1-a$） & ラット嗅皮質のスライス。$100$\,\textmu Mのアセチルコリン作動薬で、同じ細胞において、内部結合（Ib層）の電位は$60\%$以上、外部入力（Ia層）の電位は$15\%$未満低下する。抑制はシナプス前性。海馬スライス（CA1）では内部経路と入力経路で$89\%$対$40\%$ & \cite{HasselmoBower1992,HasselmoSchnell1994} & 測定済み \\
5 & 第二の作用：入力を強める（係数$\frac{1+a}{2}$） & 新皮質スライス。ニコチン性受容体は視床からの入力のシナプスを強め、皮質内のシナプスには作用しない。ムスカリン性受容体は両方を弱める。アセチルコリンはニューロンの興奮性を高める（レビュー） & \cite{Gil1997,Hasselmo2006} & 測定済み \\
6 & 第三の作用：重みの変化を容易にする（学習率$\eta a$） & ラット。基底核（皮質へのアセチルコリンの供給源）の電気刺激を音と対にすると、成体でも聴覚皮質の地図が提示音に合わせて大きく再編される & \cite{KilgardMerzenich1998,Hasselmo2006} & 測定済み \\
7 & \eqref{eq:acetnet}の第二式のまばらなパターン用ヘッブ則は、一部からパターンを補完する連想記憶を与える & まばらなパターンと共分散則をもつネットワークの容量計算。活動ニューロンの割合$p$が小さいと、ネットワークは$p=1/2$の場合よりはるかに多くのパターンを保存する & \cite{Tsodyks1988} & 理論 \\
8 & 低い$a$での書き込みは記憶を壊す：ネットワークは見たものではなく思い出したものを書き込む。$a=0.1$での破損は$a=0$より弱くない & \texttt{models/\allowbreak acet\_memory.py}：ニューロン$200$個、$20$個ずつのパターン五つ、新パターンはそのうち一つと$10$個のニューロンを共有、$10$回の実行。$a=0$では新パターンは記憶されず、古いパターンは$10$個中$5.9$個の無関係なニューロンを引き連れる。$a=0.1$では新パターンは思い出される（$0.97$）が両者が融合し、新パターンは$7.3$個、古いパターンは$7.9$個の無関係なニューロンを引き連れる。$a=0.2$からは1個未満 & 本章で導出 & モデル \\
9 & 命題~\ref{prop:writeread}：書き込みと読み出しが同じようにうまくいくレベル$a$はない（図~\ref{fig:acetsim}） & 同じスクリプト、学習率$\eta a$。新パターンが1回の提示で丸ごと記憶されるのは$a\ge0.9$、$a=0.75$で$0.8$、$a\le0.5$では記憶されない。パターンの半分からの読み出し：$a\le0.2$で$1.0$、$a=0.3$で$0.96$、$a=0.4$で$0.04$ & 本章で導出 & モデル \\
10 & 脳では1本のブロードキャスト配線、\nt{ACET}が書き込みと読み出しを切り替える & この考えはスライスでの測定に基づくモデルからとった。最も直接的な実験はヒトでのもの。ムスカリン受容体遮断薬スコポラミンは新しい単語対の暗記、特に古いものと重なる対の暗記を悪化させるが、すでに学習した対の想起は悪化させない。ネットワークの二つのモードの直接測定はない & \cite{HasselmoAndersonBower1992,Atri2004} & 仮説 \\
11 & フィルタ~\eqref{eq:kalman}は最適である。入力の重みと学習率は同じ量$P/R$で、定常状態では$P=\sqrt{qR}$、記憶は$\sqrt{R/q}$（命題~\ref{prop:kalman}） & 実験を要しない。カルマン--ビューシーフィルタの最適性は推定理論の古典的結果であり、定常状態は本章で導いた & \cite{KalmanBucy1961} & 理論 \\
12 & \nt{ACET}はネットワークに$P$に似た量、予期された不確実性を伝える & 注意の確率モデルで提案された。アセチルコリンのレベルが推定の不確実性に従うことの直接測定はない & \cite{YuDayan2005} & 仮説 \\
13 & \nt{ACET}ニューロンの一部は報酬と罰に20ミリ秒ほどで応答し、強化が予期しないものであるほど強く応答する & マウス、光遺伝学的に同定した前脳基底部のコリン作動性ニューロン。報酬と罰への応答は$18\pm3$\,msで、大きさは強化の予期しなさとともに増し、二つの核のニューロンで応答は似ている & \cite{Hangya2015} & 測定済み \\
14 & \nt{NORA}は世界の予期しない変化に気づき、\nt{ACET}はネットワークを書き込みモードに保つ & この分業はモデルで提案された。間接的には、ヒトで\nt{NORA}に関係する瞳孔径が変化と学習率を追跡する。\nt{NORA}--\nt{ACET}の組の直接の検証はない & \cite{YuDayan2005,Nassar2012} & 仮説 \\
15 & 徐波睡眠では、読み出しモードのネットワークが日中に書き込んだものを再生し、その再生が長期記憶への書き直しになる & ラット海馬のニューロン集団の記録。日中に一緒に働いた細胞は、その後の徐波睡眠でより頻繁に、同じ順序で一緒に発火する（測定済み）。書き直しについて。学習後に海馬のリップルを抑制すると空間記憶が悪化し、自発リップルを延ばすと改善する。原因が\nt{ACET}の低レベルであることは証明されていない & \cite{WilsonMcNaughton1994,Skaggs1996,Girardeau2009,FernandezRuiz2019,Hasselmo1999} & 仮説 \\
\bottomrule
\end{longtable}}
